\documentclass[12pt]{article}
\usepackage[utf8]{inputenc}
\usepackage[T1]{fontenc}
\usepackage{hyperref}
\title{Berta Language Resources for K-12 Curriculum Development}
\author{Compiled by Claude Sonnet 5.5 (claude-sonnet-5-5) for Andr\'as Kornai}
\date{}
\begin{document}
\maketitle

\section{Introduction}
Berta (ISO 639-3 wti; Glottocode bert1248) is a Nilo-Saharan language of the Benishangul-Gumuz Region of Ethiopia and neighbouring Sudan. Search results gave about 125,000 speakers in Ethiopia and about 22,000 in Sudan, with a Latin-based writing system. This compilation rests on a small number of web search results. Little educational material for Berta turned up.

\section{Primary Educational Resources}
I found no Berta primer, textbook or published mother-tongue curriculum in the search results. A SIL survey from 2016 is reported to have covered six Berta speech varieties in Ethiopia. A related trial-edition teacher's module, ``Communicative Mother Tongue Skills'', exists for Shinasha, a different language. It comes from the Benishangul-Gumuz Language Development Project, a joint effort of the regional Education Bureau and SIL Ethiopia. It is not a Berta resource, but it shows the kind of material that project produces. A search of the SIL Ethiopia catalogue and the regional Education Bureau may find Berta material.

\section{Linguistic and Reference Materials}
\begin{itemize}
\item Overview of the Berta language: \url{https://en.wikipedia.org/wiki/Berta_language}
\item Overview of the Berta languages group: \url{https://en.wikipedia.org/wiki/Berta_languages}
\item Verbix page on Berta: \url{https://docs.verbix.com/Languages/Berta}
\end{itemize}
I did not check these pages for depth. I found no Berta dictionary or reference grammar in the results.

\section{Religious and Community Resources}
\begin{itemize}
\item Joshua Project entry on Berta in Ethiopia (demographic and people-group information): \url{https://www.joshuaproject.net/index.php/people_groups/10793/ET}
\end{itemize}

\section{Digital and Audio Resources}
No audio, video or software resources for Berta turned up. The Rosetta Panglossia entry (\url{https://rosettapanglossia.longnow.org/index.php?title=Wti}) lists alternate names and basic information only.

\section{Curriculum Development Context}
The Foundation for Endangered Languages report (\url{https://www.ogmios.org/ogmios_files/172.htm}) says five local languages are officially recognized in Benishangul-Gumuz: Berta, Gumuz, Shinasha, Mao and Komo. It also says the regional government was developing these languages, and that how they fit into the school curriculum is often decided on political grounds. A Berta curriculum would probably have to be built from scratch, working with the regional Education Bureau and SIL Ethiopia.

\end{document}